\section{2025-01-23}

\starttyping
SELECT Vorname, Nachname, Klasse, 'Schuelerliste1' as Schule FROM Schuelerliste1 WHERE (...) AND ... ORDER BY Nachnamen DESC
UNION
SELECT Vorname, Nachname, Klasse, 'Schuelerliste2' as Schule FROM Schuelerliste2 WHERE (...) AND ... ORDER BY Nachnamen DESC
ORDER BY Nachnahme DESC;
\stoptyping

\subsection{S. 319 A2}

\startitemize[packed,a]
\item Die drei Ansichten können mit auf Basis folgenden Schema aufgebaut werden:

Programm(ID, Name, Fortschritt, Sender, Genre, Zeit, Art, FKS, Tags, Folge, Dauer, Erscheinungsjahr, Beschreibung, Cover)

In dieser Tabelle sind alle benötigten Informationen.

\item 

\startframedtext[width=\dimexpr\textwidth+\rightmarginwidth+\rightmargindistance\relax,frame=off,offset=none] % This spans textwidth + margin width
\startplacetable[location={force},title={Beispiel Datensatz}]
\setupTABLE[r][1][background=color,backgroundcolor=gray]
\bTABLE
\startCSV
SendungsID, Sender, Sendung, Datum, Anfang, Ende, Genre, Infotext
1, DasErste, Drei Fragezeichen, 20250101, 0000, 0020, Horror, "Die drei Fragezeichen sind auf der Scuher nach dem geheimen Artefakt, doch Justus Jonas hat andere Ideen"
2, DasZweite, Wer wird Millionär?, 20261204, 1000, 1130, Satiere, "Zehn Fragen, ein Gewinner. Wer wird dieses mal gewinnen."
3, RTL, NuNaJan, 20261104, 1130, 1230, Kinderverdummung, "keine Ahnung was hier abgeht"
4, Nikelodion, Spongebob, 20260804, 1230, 1330, Kinderverdummung, "Ein Schwamm der unter wasser Lebt"
\stopCSV
\eTABLE
\stopplacetable
\stopframedtext

\item Die SQL-Abfragen
\startitemize[packed,r]
\item Alle Sender, die die Tagesschau ausstrahlen
\item Alle Nachrichten, die nach 10 Uhr ausgestrahlt werden, mit Sender, Sendung und wann sie beginnen
\item Alle Sender die "schau" im Namen beinhlaten
\item Es werden angegeben, wie viele Sendungen in einem Genre existieren, die nicht zu ZDFHD als Sender haben. Die Genre Show und Sport werden ausgeschlossen.
\item Listet alle einzelnen Genres auf
\item Alle Sendungen, die nach der Tagesschau ausgestrahlt werden.
\stopitemize
\stopitemize


